% Esquisse d'un programme (cote 119, copy 2): archive cote 119, site page 87 (PDF page 88), typescript page 9.
% typescript. Draft compared: esquisse/pages/esq_119_011.tex (copy 1, site 11).
% AI-transcribed by Claude Opus 5.5 (claude-opus-5-5, adaptive thinking, effort high), Message Batches API, batch esq119_full (msgbatch_015xqDHkTumDysRF1qngZhRJ), 2026-11-13; not yet checked by eye.
\scanpage{119/87}
\mspage{9}

permettent de considérer ces surfaces comme associées à des courbes
algébriques complexes (et même sur $\overline{\mathbb{Q}}$) canoniquement déterminées.
Il y a là un jeu de chassé-croisé typique entre le combinatoire,
et l'algébrique complexe (ou mieux, l'algébrique sur $\overline{\mathbb{Q}}$).

La ``petite boite'' aux pièces toutes identiques, qui a le
charme de l'économie, donnera sans doute une description relative-%
ment compliquée pour les relations (compliquée, mais nullement inex-%
tricable !). La grande boite donnera lieu à des relations plus nom-%
breuses (du fait qu'il y a beaucoup plus de points-bases et de che-%
mins remarquables entre eux), mais à structure plus transparente. Je
prévois qu'en dimension modulaire 2, tout comme dans le cas plus ou
moins familier de la dimension modulaire 1 (avec notamment la des-%
\margin{//} cription de $S\ell(2,\mathbb{Z})$ par $(\rho,\sigma \mid \rho^3 = \sigma^2,\ \over{\rho}{\sigma}^4 = \over{\sigma}{\rho}^6 = 1)$), on trouvera
un engendrement par les groupes d'automorphismes des trois types de
pièces pertinentes, avec des relations simples que je n'ai pas déga-%
gé à l'heure d'écrire ces lignes. Peut-être même trouvera-t-on un
principe de ce genre pour tous les $T_{g,\nu}$, ainsi qu'une décomposition
cellulaire de $M^{\wedge}_{g,\nu}$ généralisant celles qui se présentent sponta-%
nément pour $M^{\wedge}_{0,4}$ et $M^{\wedge}_{1,1}$, et que j'entrevois dès à présent pour
la dimension modulaire 2, en utilisant les hypersurfaces correspon-%
dant aux diverses \uline{structures réelles} sur les structures complexes
envisagées, pour effectuer le découpage cellulaire voulu.


3. Plutôt que de suivre (comme prévu) un ordre thématique ri-%
goureux, je me suis laissé emporter par ma prédilection pour un thème
particulièrement riche et brûlant, auquel je compte me consacrer
d'ailleurs prioritairement pendant quelque temps, à partir de la ren-%
trée 84/85. Je reprends donc l'exposé thématique là où je l'ai laissé,
tout au début du paragraphe précédent.

Mon intérêt pour les surfaces topologiques commence à poindre
en 1974, où je propose à Yv\unsure{è|e}s Ladegaillerie le thème de l'étude iso-%
topique des plongements d'un 1-complexe topologique dans une surface
compacte. Dans les deux années qui suivent, cette étude le conduit à
un remarquable théorème d'isotopie, donnant une description algébrique

\margin{87}
